\chapter{Periods, Torelli, and lattice polarizations}\label{ch:torelli}

The theme of this chapter is a single sentence: \emph{a K3 surface is
determined by the position of the line $H^{2,0}$ inside the fixed lattice
$\Lambda_{K3}$, and every allowable position occurs.} Its two halves are the
Torelli theorem and the surjectivity of the period map. Dolgachev's theory of
lattice-polarized K3 surfaces is then the systematic way to cut out the
families one actually studies, and it produces the rank-3 lattice
$T_n=\U\oplus\lat{2n}$ of Chapter~\ref{ch:lattices} as the transcendental
lattice of a one-parameter family whose parameter space is a modular curve.

\section{The period domain}

Fix a \emph{marking}, an isometry $\phi\colon H^2(X,\Z)\to\Lambda:=\Lambda_{K3}$.
The line $\phi(H^{2,0}(X))=\C\sigma$ satisfies the Riemann bilinear relations
\[
  \langle\sigma,\sigma\rangle=0,\qquad\langle\sigma,\bar\sigma\rangle>0 ,
\]
so it is a point of the \emph{period domain}
\[
  \Omega:=\{[\sigma]\in\PP(\Lambda_\C):\sigma^2=0,\ \sigma\bar\sigma>0\},
\]
an open subset of a $20$-dimensional quadric, with two connected components
exchanged by conjugation. The holomorphic $2$-form is the only datum in
$H^2$ that varies with the surface.

\begin{theorem}[\Lstatus, Local Torelli]
The period map from the base of a local universal deformation of $X$ to
$\Omega$ is a local isomorphism. In particular K3 surfaces have $20$ moduli
\cite[Ch.~6, Prop.~2.8]{Huybrechts2016}.
\end{theorem}

\begin{theorem}[\Lstatus, Global Torelli]\label{thm:torelli}
Two K3 surfaces $X,X'$ are isomorphic if and only if there is a Hodge isometry
$H^2(X,\Z)\cong H^2(X',\Z)$. Moreover, a Hodge isometry that maps an ample
(or K\"ahler) class to an ample (K\"ahler) class is induced by a unique
isomorphism $X'\to X$
\cite{PiatetskiShapiroShafarevich1971},~\cite{BurnsRapoport1975},
\cite[Thm.~7.5.3]{Huybrechts2016}.
\end{theorem}

\begin{theorem}[\Lstatus, surjectivity of the period map]\label{thm:surj}
Every point of $\Omega$ is the period of some marked K3 surface
\cite{Todorov1980,Looijenga1981}, \cite[Thm.~7.4.1]{Huybrechts2016}.
\end{theorem}

Torelli is the statement that the lattice remembers everything; surjectivity,
that the lattice forgets nothing. Together they turn questions about K3
surfaces into questions about lines in $\Lambda_\C$ and sublattices of
$\Lambda$ --- which is why Chapter~\ref{ch:lattices} came first.

\begin{example}[Picard number from the period]
Given $[\sigma]\in\Omega$, the N\'eron--Severi lattice of the corresponding
surface is $\sigma^{\perp}\cap\Lambda$ (Lefschetz $(1,1)$). A generic
$[\sigma]$ has $\sigma^{\perp}\cap\Lambda=0$, so the generic K3 surface has
$\rho=0$ and is not projective. Requiring $\rho\ge1$ is one condition;
requiring $\rho\ge19$ is nineteen, leaving a one-dimensional family; and
$\rho=20$ is a countable set of points --- the singular K3 surfaces of
Theorem~\ref{thm:shioda-inose}, with $T=\sigma^{\perp\perp}\cap\Lambda$ of rank~$2$.
\end{example}

\section{Lattice-polarized K3 surfaces}

\begin{definition}[Dolgachev~\cite{Dolgachev1996}]
Let $M$ be an even hyperbolic lattice (signature $(1,\rk M-1)$) admitting a
primitive embedding $M\hookrightarrow\Lambda$. An \emph{$M$-polarized K3
surface} is a pair $(X,j)$ with $j\colon M\hookrightarrow\NS(X)$ a primitive
embedding such that $j(M)$ contains a pseudo-ample class (one in the closure
of the ample cone, i.e.\ nef and big). Two such pairs are isomorphic if there
is an isomorphism of surfaces compatible with the embeddings.
\end{definition}

Fix a primitive embedding $M\subset\Lambda$ and let $T:=M^{\perp}\subset\Lambda$,
of signature $(2,20-\rk M)$... more precisely $(2,\,20-\rk M+1-1)=(2,19-\rk M+1)$;
for $\rk M=19$ this is $(2,1)$. The period of an $M$-polarized surface lies in
\[
  \Omega_M:=\{[\sigma]\in\Omega:\sigma\perp M\}
  =\{[\sigma]\in\PP(T_\C):\sigma^2=0,\ \sigma\bar\sigma>0\},
\]
a domain of dimension $\rk T-2=20-\rk M$.

\begin{theorem}[\Lstatus, Dolgachev]\label{thm:dolgachev-moduli}
The coarse moduli space of $M$-polarized K3 surfaces is
$K_M=\Gamma_M\backslash\Omega_M$ where $\Gamma_M$ is the group of isometries of
$\Lambda$ acting trivially on $M$ (identified with a subgroup of $\Orth(T)$).
Its dimension is $20-\rk M$, and the period map is an isomorphism onto the
complement of a locally finite union of divisors (the ``$(-2)$-walls'')
\cite[Thm.~3.1, Prop.~3.3, Thm.~7.3]{Dolgachev1996}.
\end{theorem}

The walls are where the pseudo-ample condition fails: the period becomes
orthogonal to a $(-2)$-vector $\delta\in T$, which is then an algebraic class
not in $M$ but orthogonal to it, and the surface acquires a rational curve
that the polarization did not see. We shall meet exactly two such walls in the
family of Chapter~\ref{ch:frontier}, and Chapter~\ref{ch:lattices}'s
Examples~\ref{ex:fricke-root} and~\ref{ex:root7} are their roots.

\section{The $M_n$-polarized families}

Take $M=M_n=\U\oplus E_8(-1)^2\oplus\lat{-2n}$, rank $19$, so $T=T_n=\U\oplus\lat{2n}$
(Example~\ref{ex:Mn}) and $K_{M_n}$ is a curve.

\begin{theorem}[\Lstatus, Dolgachev~\cite{Dolgachev1996}, \S7]\label{thm:Mn}
\begin{enumerate}
\item $(M_n)^{\perp}=\U\oplus\lat{2n}$, uniquely.
\item $\Orth^{+}(\U\oplus\lat{2n})/\{\pm1\}\cong\Gamma_0(n)^{+}$, the extension
  of $\Gamma_0(n)$ by the Fricke involution, and
  $K_{M_n}\cong\Gamma_0(n)^{+}\backslash\HH$, the Fricke modular curve
  $X_0(n)^{+}$ (Chapter~\ref{ch:modular}).
\item The $(-2)$-walls removed from the ample locus form an explicit countable
  set (Thm.~7.3 there).
\item $K_{M_n}$ has a canonical coordinate, an integer-coefficient Hauptmodul
  when $X_0(n)^+$ has genus $0$ (Thm.~7.5 there).
\end{enumerate}
\end{theorem}

The identification in (ii) can be made completely explicit, and it is
instructive to see the isometry group of a rank-3 lattice \emph{be} a modular
group. Write the period as
\[
  \omega(\tau)=e-n\tau^2f+\tau w\in T_n\otimes\C,\qquad
  \omega(\tau)^2 = 2\cdot1\cdot(-n\tau^2)+2n\tau^2=0 ,
\]
so $\tau\mapsto[\omega(\tau)]$ maps $\HH$ onto a component of $\Omega_{M_n}$
(the isotropic vectors of the rank-3 lattice form a conic, and $\tau$ is its
rational parameter). Then one checks:

\begin{proposition}[\Kstatus\ for the identities, \Lstatus\ for the group-theoretic converse]\label{prop:rho}
There is an explicit representation $\rho\colon\SL_2(\Z)\to\GL_3(\Z)$
(a symmetric square in the basis $(e,f,w)$) such that for
$\gamma=\bigl(\begin{smallmatrix}a&b\\nc&d\end{smallmatrix}\bigr)\in\Gamma_0(n)$,
\[
  \rho(\gamma)^{\mathsf T}\,T_n\,\rho(\gamma)=T_n,\qquad
  \rho(\gamma)\,\omega(\tau)=(nc\tau+d)^2\,\omega(\gamma\tau),
\]
the second identity holding as a polynomial identity in $\tau$ with no
determinant hypothesis; the swap $\sigma$ of Example~\ref{ex:swap} satisfies
$\sigma\,\omega(\tau)=(-n\tau^2)\,\omega(-1/(n\tau))$; and the Atkin--Lehner
elements $W_Q$ for $Q\parallel n$ act by integer matrices in the same way.
That these exhaust $\Orth^+(T_n)/\pm1$ is Dolgachev's theorem, quoted.
\end{proposition}

The kernel-checked identities are in \artifact{Agora/Geometry/ModularAction.lean}
(\artifact{rho\_isometry}, \artifact{rho\_mul}, \artifact{rho\_mulVec\_period},
\artifact{rhoQ\_*}) and \artifact{MnLattice.lean} (\artifact{swap\_is\_fricke});
see Chapter~\ref{ch:formal}. The reader should notice what they do and do not
say: they exhibit the modular group \emph{inside} the isometry group, with the
automorphy factor $(nc\tau+d)^2$ of weight $2$ appearing as the scalar by
which a lattice isometry rescales the period. The equality of the two groups
is not among them.

\begin{remark}[the two rank-3 lattices]\label{rem:two-lattices}
There is a second rank-3 lattice on which $\Gamma_0(n)$ acts by symmetric
squares: the lattice of binary quadratic forms $ax^2+bxy+cy^2$ with $n\mid a$,
with the discriminant form $b^2-4nac$ (up to normalization), of discriminant
$-4n^2$. It is \emph{not} isometric to $T_n$, whose discriminant is $-2n$
(an isometry preserves the discriminant; $-4n^2=-2n$ forces $n\in\{0,\tfrac12\}$).
Both carry $\Gamma_0(n)^+$-actions, both have signature $(2,1)$, and they have
been conflated in at least one working document. Keep them apart: the
transcendental lattice is $T_n$; the form lattice is a different object that
happens to share a group.
\end{remark}

\section{Exercises}

\begin{exercise}
Show that $\Omega$ has exactly two connected components and that they are
interchanged by $\sigma\mapsto\bar\sigma$. Show that for $T$ of signature
$(2,1)$, $\Omega_M$ is two copies of the upper half-plane.
\end{exercise}

\begin{exercise}[$\star$]
Verify by hand that $\omega(\tau)=e-n\tau^2f+\tau w$ is isotropic in $T_n$,
that $\langle e-f,\omega(\tau)\rangle=-(n\tau^2+1)$, and hence that the wall
of the Fricke root $e-f$ is the locus $n\tau^2=-1$, i.e.\ $\tau=i/\sqrt n$ in
$\HH$. Compute the wall of the root $(2,-4,1)$ in $T_7$ and show it is
$\tau=\tfrac12\pm i/\sqrt{28}$; compute the wall of $(-2,4,1)$ and show its
roots differ from these by $1$.
\end{exercise}

\begin{exercise}
Write down $\rho(\gamma)$ for
$\gamma=\bigl(\begin{smallmatrix}1&1\\0&1\end{smallmatrix}\bigr)$ (translation)
and check the period identity directly.
\end{exercise}

\section*{Chapter note}

Theorems~\ref{thm:torelli}--\ref{thm:surj} and the lattice-polarized theory
are quoted from Huybrechts~\cite{Huybrechts2016} and Dolgachev~\cite{Dolgachev1996};
Dolgachev's \S7 was read from a hash-pinned copy in the program's literature
manifest. Proposition~\ref{prop:rho} states what the kernel checked, as
transcribed from the repository at the commit of Chapter~\ref{ch:formal}.
